% GENERATED FILE. Edit canonical lessons, then run npm run generate:books.
% canonical-source-hash: fnv1a64:fa061ca7bb284f4c
% canonical-lessons: ZH-C469-chongdong, ZH-C469-wenhe, ZH-W469-wen-warm, ZH-C469-zhengjing, ZH-C469-wenzhong, ZH-C469-jixingzi, ZH-R469-review-464-469

\chapter{Gentle and Steady}
\label{ch:zh-gentle-and-steady}
\hlchaptermodality{469}

\begin{chapteropening}
\textbf{By the end of this chapter:} \emph{I can say someone is impulsive, gentle, decent or steady, and that someone is an impatient sort.}

Complete the last lesson of chapter 469: I can say someone is impulsive, gentle, decent or steady, and that someone is an impatient sort.
\end{chapteropening}

\section[chōngdòng]{chōngdòng — impulsive; rash}

\label{lesson:ZH-C469-chongdong}



\begin{quote}
Before the new one: say the Mandarin for quick to pitch in, then the Mandarin for thoughtful of others' needs.
\end{quote}

\subsection*{You'll want to know: chōngdòng}
\emph{chōngdòng} — \textquotedblleft{}impulsive; rash\textquotedblright{}.

\begin{quote}
\emph{tā tài chōngdòng le chángcháng shuō cuò huà} — He's too impulsive and often says the wrong thing.
\end{quote}

Its characters come later; for now it is a word for the ear.

\begin{cousinweb}
\emph{chōngdòng} is built from \emph{chōng} 'to dash against' and \emph{dòng} 'to move', both inherited from Middle Chinese, \emph{tsyhowng} and \emph{duwngX} (Guangyun, in Baxter's transcription), which glosses the first 'to rush at' and the second 'restless' and 'to move'.
\end{cousinweb}

\begin{grammarlens}[title={What you've built}]
One more word for saying how it was.
\end{grammarlens}

\subsection*{Your turn}
\begin{itemize}
\raggedright
  \item \emph{Say it:} \emph{chōngdòng}
  \item \emph{Say it:} \emph{chōngdòng}, once more
  \item \emph{Say it:} \emph{tā tài chōngdòng le chángcháng shuō cuò huà} — He's too impulsive and often says the wrong thing.
  \item \emph{Recall:} say the Mandarin for quick to pitch in, then the Mandarin for thoughtful of others' needs, then say \emph{chōngdòng} again
\end{itemize}

\begin{tcolorbox}[breakable,colback=teal!4,colframe=teal!35!black,title={Before you move on}]
What does chōngdòng mean? (\textquotedblleft{}Impulsive; rash\textquotedblright{}.) Say it once more.
\end{tcolorbox}
\section[\texorpdfstring{\zh{温和} (wēnhé)}{温和 (wēnhé)}]{\zh{温和} (wēnhé) — gentle in manner; soft-spoken}

\label{lesson:ZH-C469-wenhe}



\begin{quote}
Before the new one: say the Mandarin for thoughtful of others' needs, then the Mandarin for impulsive.
\end{quote}

\subsection*{You'll want to know: \zh{温和}}
\textbf{\zh{温和}} — \emph{wēnhé} — \textquotedblleft{}gentle in manner; soft-spoken\textquotedblright{}.

\begin{quote}
\textbf{\zh{他的个性比我温和}} \emph{tā de gèxìng bǐ wǒ wēnhé} — He has a gentler nature than I do.
\end{quote}

The new character, in pieces: \zh{温} — twelve strokes in three pieces: on the left the water piece, a dot, a second dot and a rising stroke; at the top right \zh{日} \emph{rì}, four strokes; at the foot a dish of five strokes, a vertical, a stroke that runs right and turns down, two short verticals and a bar. The next lesson puts it in your hand.

\begin{cousinweb}
\textbf{\zh{温和}} (\emph{wēnhé}) is built from \zh{温} 'warm' and \zh{和} 'harmony', both inherited from Middle Chinese, \emph{'won} and \emph{hwa} (Guangyun, in Baxter's transcription), which glosses the first 'mild', 'good', 'soft' and 'warm' and the second 'in tune'.
\end{cousinweb}

\begin{grammarlens}[title={What you've built}]
One more word for saying how it was.
\end{grammarlens}

\subsection*{Your turn}
\begin{itemize}
\raggedright
  \item \emph{Say it:} \emph{wēnhé}
  \item \emph{Say it:} \emph{wēnhé}, once more
  \item \emph{Say it:} \emph{tā de gèxìng bǐ wǒ wēnhé} — He has a gentler nature than I do.
  \item \emph{Recall:} say the Mandarin for thoughtful of others' needs, then the Mandarin for impulsive, then say \emph{wēnhé} again
\end{itemize}

\begin{tcolorbox}[breakable,colback=teal!4,colframe=teal!35!black,title={Before you move on}]
What does \zh{温和} mean? (\textquotedblleft{}Gentle in manner; soft-spoken\textquotedblright{}.) Say it once more.
\end{tcolorbox}
\section[\texorpdfstring{\zh{温} (wēn)}{温 (wēn)}]{\zh{温} — the water piece, then \zh{日} over a dish}

\label{lesson:ZH-W469-wen-warm}



\begin{quote}
The word first, then the character inside it.

\begin{itemize}
\raggedright
  \item \emph{Recall:} say \emph{wēnhé} — \textbf{R1}, one lesson back, before it decays
\end{itemize}

\begin{quote}
\zh{汉} \zh{日} \zh{和}
\end{quote}

\emph{Read it:} each one aloud, with nothing said first — \emph{hàn}, \emph{rì}, \emph{hé}
\end{quote}

\subsection*{Script — the water piece, then \zh{日} over a dish}
\hlblockfigure{\detokenize{figures/ZH-W469-wen-warm-filmstrip.pdf}}{How \zh{温} is written, stroke by stroke}

\begin{quote}
\zh{温}
\end{quote}

\textbf{Three pieces: the left, then the top right, then the foot.} Count the pieces, not the strokes.

\textbf{The left piece, three strokes: the water piece, as in \zh{汉}.} A dot, a second dot under it, then a stroke that rises from the bottom left up to the right.

\textbf{The top right, \zh{日}, four strokes.} A vertical down the left, a stroke that runs right and turns down, a bar inside, then the bar across the foot.

\textbf{The foot, a dish, five strokes.} A short vertical on the left, a stroke that runs right and turns down, two short verticals inside, then a long bar across the bottom, under everything.

3 + 4 + 5 = twelve.

\subsection*{Your turn}
\begin{itemize}
\raggedright
  \item \emph{Write it:} \zh{温}, all twelve, with the model in front of you
  \item \emph{Cover:} write it again from memory, and count the pieces, not the strokes
\end{itemize}

\begin{tcolorbox}[breakable,colback=teal!4,colframe=teal!35!black,title={Before you move on}]
How many pieces does \zh{温} have? (\textbf{Three: three strokes, four and five}.) Which written character sits at the top right? (\textbf{\zh{日}}.)

Source: \href{https://raw.githubusercontent.com/chanind/hanzi-writer-data/68d10a4b21150cae5e1ebbd223eed289cf32d90c/data/\%E6\%B8\%A9.json}{Hanzi Writer Data \zh{温}, pinned stroke paths and medians}.
\end{tcolorbox}
\section[\texorpdfstring{\zh{正经} (zhèngjīng)}{正经 (zhèngjīng)}]{\zh{正经} (zhèngjīng) — decent; serious and proper}

\label{lesson:ZH-C469-zhengjing}



\begin{quote}
Before the new one: say the Mandarin for impulsive, then the Mandarin for gentle in manner.
\end{quote}

\subsection*{You'll want to know: \zh{正经}}
\textbf{\zh{正经}} — \emph{zhèngjīng} — \textquotedblleft{}decent; serious and proper\textquotedblright{}.

\begin{quote}
\textbf{\zh{他是一个很正经的人}} \emph{tā shì yí ge hěn zhèngjīng de rén} — He's a very decent man.
\end{quote}

\begin{cousinweb}
\textbf{\zh{正经}} (\emph{zhèngjīng}) is built from \zh{正} 'straight' and \zh{经} 'the warp thread', both inherited from Middle Chinese, \emph{tsyengH} and \emph{keng} (Guangyun, in Baxter's transcription), which glosses the first 'right' and 'proper' and the second 'constant' and as the warp of cloth. The sense 'decent' is not given there.
\end{cousinweb}

\begin{grammarlens}[title={What you've built}]
One more word for saying how it was.
\end{grammarlens}

\subsection*{Your turn}
\begin{itemize}
\raggedright
  \item \emph{Say it:} \emph{zhèngjīng}
  \item \emph{Say it:} \emph{zhèngjīng}, once more
  \item \emph{Say it:} \emph{tā shì yí ge hěn zhèngjīng de rén} — He's a very decent man.
  \item \emph{Recall:} say the Mandarin for impulsive, then the Mandarin for gentle in manner, then say \emph{zhèngjīng} again
\end{itemize}

\begin{tcolorbox}[breakable,colback=teal!4,colframe=teal!35!black,title={Before you move on}]
What does \zh{正经} mean? (\textquotedblleft{}Decent; serious and proper\textquotedblright{}.) Say it once more.
\end{tcolorbox}
\section[wěnzhòng]{wěnzhòng — steady; calm and reliable}

\label{lesson:ZH-C469-wenzhong}



\begin{quote}
Before the new one: say the Mandarin for gentle in manner, then the Mandarin for decent and proper.
\end{quote}

\subsection*{You'll want to know: wěnzhòng}
\emph{wěnzhòng} — \textquotedblleft{}steady; calm and reliable\textquotedblright{}.

\begin{quote}
\emph{tā bǐ tā jiějie wěnzhòng} — She's steadier than her big sister.
\end{quote}

Its characters come later; for now it is a word for the ear.

\begin{cousinweb}
\emph{wěnzhòng} is built from \emph{wěn} 'settled' and \zh{重} 'heavy', both inherited from Middle Chinese, \emph{'wonX} and \emph{drjowngX} (Guangyun, in Baxter's transcription), which glosses the first as grain gathered and 'safe and settled', and the second 'thick' and 'careful'.
\end{cousinweb}

\begin{grammarlens}[title={What you've built}]
One more word for saying how it was.
\end{grammarlens}

\subsection*{Your turn}
\begin{itemize}
\raggedright
  \item \emph{Say it:} \emph{wěnzhòng}
  \item \emph{Say it:} \emph{wěnzhòng}, once more
  \item \emph{Say it:} \emph{tā bǐ tā jiějie wěnzhòng} — She's steadier than her big sister.
  \item \emph{Recall:} say the Mandarin for gentle in manner, then the Mandarin for decent and proper, then say \emph{wěnzhòng} again
\end{itemize}

\begin{tcolorbox}[breakable,colback=teal!4,colframe=teal!35!black,title={Before you move on}]
What does wěnzhòng mean? (\textquotedblleft{}Steady; calm and reliable\textquotedblright{}.) Say it once more.
\end{tcolorbox}
\section[jíxìngzi]{jíxìngzi — an impatient sort; a hothead}

\label{lesson:ZH-C469-jixingzi}



\begin{quote}
Before the new one: say the Mandarin for decent and proper, then the Mandarin for steady and reliable.
\end{quote}

\subsection*{You'll want to know: jíxìngzi}
\emph{jíxìngzi} — \textquotedblleft{}an impatient sort; a hothead\textquotedblright{}.

\begin{quote}
\emph{wǒ bàba shì ge jíxìngzi} — My dad is the impatient sort.
\end{quote}

Its characters come later; for now it is a word for the ear.

\begin{cousinweb}
\emph{jíxìngzi} is built from \emph{jí} 'hasty', \zh{性} 'nature' and \zh{子}, a noun ending, all inherited from Middle Chinese, \emph{kip}, \emph{sjengH} and \emph{tsiX} (Guangyun, in Baxter's transcription), which glosses the first 'urgent' and 'quick'.
\end{cousinweb}

\begin{grammarlens}[title={What you've built}]
One more word for saying how it was.
\end{grammarlens}

\subsection*{Your turn}
\begin{itemize}
\raggedright
  \item \emph{Say it:} \emph{jíxìngzi}
  \item \emph{Say it:} \emph{jíxìngzi}, once more
  \item \emph{Say it:} \emph{wǒ bàba shì ge jíxìngzi} — My dad is the impatient sort.
  \item \emph{Recall:} say the Mandarin for decent and proper, then the Mandarin for steady and reliable, then say \emph{jíxìngzi} again
\end{itemize}

\begin{tcolorbox}[breakable,colback=teal!4,colframe=teal!35!black,title={Before you move on}]
What does jíxìngzi mean? (\textquotedblleft{}An impatient sort; a hothead\textquotedblright{}.) Say it once more.
\end{tcolorbox}
\section[(dialogue)]{Looking back: from Losing My Temper to Gentle and Steady}

\label{lesson:ZH-R469-review-464-469}



\begin{quote}
Say the Mandarin for to lose one's temper and a hot temper.
\end{quote}

\subsection*{Your turn}
\begin{itemize}
\raggedright
  \item \emph{Say it:} to lose one's temper, a hot temper, to be infuriating, to take it out on someone, furious
  \item \emph{Say it:} to one's heart's content, the fun spoiled, optimistic, pessimistic, sure of oneself
  \item \emph{Say it:} to be lost in a daze, cool-headed, a heartbeat, to give a gentle smile, peaceful inside
  \item \emph{Say it:} a person's temperament, a strong personality, an inborn trait, cheerful and outgoing, full of life
  \item \emph{Say it:} kind-hearted, pleasant to deal with, strong in spirit, quick to pitch in, thoughtful of others' needs
  \item \emph{Say it:} impulsive, gentle in manner, decent and proper, steady and reliable, an impatient sort
  \item \emph{Read it:} \zh{兴} \zh{静} \zh{性} \zh{温} — each one aloud, then a word that holds it
\end{itemize}

\begin{tcolorbox}[breakable,colback=teal!4,colframe=teal!35!black,title={Before you move on}]
To lose one's temper? (\textbf{\zh{发火}}, \emph{fāhuǒ}.) A hot temper? (\textbf{\zh{火气}}, \emph{huǒqì}.) To be infuriating? (\textbf{\zh{气人}}, \emph{qìrén}.) To take it out on someone? (\textbf{\zh{出气}}, \emph{chūqì}.) Furious? (\emph{fènnù}.) To one's heart's content? (\emph{tòngkuài}.) The fun spoiled? (\textbf{\zh{扫兴}}, \emph{sǎoxìng}.) Optimistic? (\emph{lèguān}.) Pessimistic? (\emph{bēiguān}.) Sure of oneself? (\textbf{\zh{自信}}, \emph{zìxìn}.) To be lost in a daze? (\emph{fādāi}.) Cool-headed? (\textbf{\zh{冷静}}, \emph{lěngjìng}.) A heartbeat? (\emph{xīntiào}.) To give a gentle smile? (\emph{wēixiào}.) Peaceful inside? (\emph{píngjìng}.) A person's temperament? (\emph{píqi}.) A strong personality? (\textbf{\zh{个性}}, \emph{gèxìng}.) An inborn trait? (\textbf{\zh{天性}}, \emph{tiānxìng}.) Cheerful and outgoing? (\emph{kāilǎng}.) Full of life? (\emph{huópō}.) Kind-hearted? (\emph{shànliáng}.) Pleasant to deal with? (\textbf{\zh{和气}}, \emph{héqi}.) Strong in spirit? (\emph{jiānqiáng}.) Quick to pitch in? (\emph{qínkuài}.) Thoughtful of others' needs? (\emph{tǐtiē}.) Impulsive? (\emph{chōngdòng}.) Gentle in manner? (\textbf{\zh{温和}}, \emph{wēnhé}.) Decent and proper? (\textbf{\zh{正经}}, \emph{zhèngjīng}.) Steady and reliable? (\emph{wěnzhòng}.) An impatient sort? (\emph{jíxìngzi}.)
\end{tcolorbox}
